\section{Conclusion}
\label{sec:conclusion}

GeoScope evaluates a single scientific requirement across two complementary atomic source collections: whether a language model's selected measurement and mathematical operation physically justify its conclusion. In the 55-system UMA collection, full-definition checking rejected 81 additional invalid Qwen selections relative to a comparison-ID check without excluding any of the 312 licensed proofs. The geometry analysis explains why this matters: whole-molecule displacement and internal deformation are distinct quantities even when calculated from the same structures. Correct final verdicts, therefore, cannot substitute for proof validity.

Using 48 OC20 DFT trajectory systems and four models, we found that the same class of evidence-selection failures persisted under a separate source and evaluation protocol. After removing only enclosing Markdown fences, the strongest tested model produced 567 valid proofs from 960 answerable questions (59.1\%). No model gave a qualifying abstention on its 576 insufficient-evidence questions. Model confidence could rank some valid proofs, but 185 of 601 Qwen3-32B selections with decision-token scores above 0.99 were still invalid. These tests establish complementary limits of evidence selection, operation choice, output formatting, and token likelihood; their differing scoring policies do not constitute identical replications or a controlled model-scaling experiment.

The central contribution is a physically defined, inspectable test of scientific justification rather than a new catalyst, a general scientific-reasoning solution, or a superior hybrid verifier. For autonomous materials research, it motivates keeping measurement definitions attached to claims and separating an invalid proof from genuinely missing evidence. Prospective evidence recovery, expert chemical review, and closed-loop validation are needed before improved materials-discovery outcomes can be claimed.
